\documentclass[10pt,a4paper]{letter}
\usepackage[T1]{fontenc}
\usepackage[margin=2.0cm]{geometry}
\usepackage[hidelinks]{hyperref}
\setlength{\longindentation}{0.5\textwidth}
\signature{Aleksey Buyanov\\Independent Researcher\\\href{mailto:pppuu7@gmail.com}{pppuu7@gmail.com}\\ORCID 0009-0001-2621-9305}
\address{Aleksey Buyanov\\Independent Researcher, Moscow, Russia}
\begin{document}
\begin{letter}{Editors, \textit{Quantum Science and Technology}\\IOP Publishing}
\opening{Dear Editors,}

Please consider ``Nuisance-aware resource closure and traceable calibration for quantum sensing: Raman interferometry with a Ramsey-clock transfer benchmark'' for publication as a \textbf{Paper} in \textit{Quantum Science and Technology}.

The manuscript addresses a general experimental-design problem in quantum sensing and metrology: strong raw sensitivity does not guarantee that a target parameter remains estimable when calibration, drift, readout, or transfer-function nuisance directions are inferred from the same experiment. We derive a nuisance-profiled resource law in which allocation is optimized only after shared nuisance geometry is included. The resulting information map is monotone and concave in additive resource and yields explicit maximum-information, minimum-budget, and marginal information-per-cost rules. It therefore distinguishes a true shortage of exposure from structural non-identifiability that cannot be cured by collecting more of the same data.

The main apparatus-level demonstration is a source-grounded three-pulse $^{87}$Rb Raman atom gravimeter. The calculation uses published timing, phase noise, detector noise, and vibration-spectrum information; propagates a source-traceable colored covariance; works with the transition-probability likelihood; profiles phase, drift, contrast, readout gain, common scale, and finite-reference nuisances jointly; and closes the absolute scale with a prospective differential Raman-chirp reference. Required negative controls remain non-identifiable, while the declared modulated science direction remains locally estimable.

To test whether the criterion is specific to atom gravimetry, we additionally apply the identical frozen profiler to a reduced interleaved Ramsey/atomic-clock model with local-oscillator offset and drift nuisances. The benchmark reproduces nuisance-induced failure, recovery through response diversity, continuous recovery through explicit calibration, and resource-allocation dependence without modifying the decision rule. A deterministic robustness sweep confirms the calibration law and the monotonic allocation degradation. These normalized clock results are deliberately presented as a cross-architecture design-principle test, not as a clock-performance forecast.

We believe this combination fits QST's quantum-sensing and quantum-metrology scope because the central contribution is a reusable criterion for deciding whether additional exposure, response diversity, or traceable calibration is the resource that actually limits inference. The manuscript does not claim an experimental departure from established physics and does not depend on a particular microscopic gravity model. The full computational chain is openly documented and independently reproduced.

The manuscript includes a data-availability statement and disclosure of generative-AI assistance, and is prepared as a single-anonymous review PDF with embedded figures and tables.

Thank you for your consideration.

\closing{Sincerely,}
\end{letter}
\end{document}
